\documentclass[11pt, a4paper]{letter}
\usepackage[utf8]{inputenc}
\usepackage[spanish,es-noshorthands]{babel}
\usepackage[margin=2cm]{geometry}
\usepackage{setspace}

\signature{Prof. César Rodríguez, C.I: 5701716 \\ Coordinador del Grupo GIDEAL}
\address{Universidad de Oriente \\ Vicerrectorado Académico \\ Consejo de Investigación}
\date{julio 2026}

\begin{document}
\begin{letter}{Ciudadano(a)\\Presidente(a) y demás miembros del\\Consejo de Investigación de la Universidad de Oriente (CIUDO)\\Presente.-}

\opening{Estimados señores:}

Tengo el agrado de dirigirme a ustedes en la oportunidad de presentar formalmente la inscripción del \textbf{Grupo de Investigación en Diseño Electrónico y Aprendizaje con LLMs (GIDEAL)}, un equipo de investigación multidisciplinario adscrito a esta prestigiosa casa de estudios.

Nuestro grupo tiene como propósito fundamental la convergencia entre la Ingeniería Electrónica, el Internet de las Cosas (IoT) y la Inteligencia Artificial. Con una visión innovadora enfocada en el desarrollo académico y técnico, diseñamos y desplegamos sistemas automatizados basados en agentes de IA y Grandes Modelos de Lenguaje (LLMs). 

Entre nuestras principales líneas de investigación y actividades destacan:
\begin{itemize}
    \item La implementación de sistemas conversacionales y chatbots académicos (RAG) para optimizar el procesamiento y consulta de material instruccional.
    \item La orquestación de flujos de trabajo multimodales para realizar una evaluación preliminar automatizada (sujeta a revisión del profesor) de prácticas de laboratorio y exámenes.
    \item El desarrollo de agentes inteligentes capaces de interpretar y extraer listas de materiales (BOM) y netlists desde esquemas de circuitos, conectando formatos académicos (LaTeX/circuitikz) con herramientas profesionales de automatización de diseño (EDA).
    \item La integración de servidores locales (MCP) y pasarelas de comunicación para modernizar y facilitar el diagnóstico y administración en los laboratorios.
\end{itemize}

A través del grupo GIDEAL, estamos comprometidos con la generación de soluciones tecnológicas que eleven la calidad de la enseñanza, promuevan la investigación de vanguardia en electrónica y consoliden a nuestra Universidad en la frontera de la inteligencia computacional aplicada. Anexo a esta comunicación se presentan los recaudos, planillas y credenciales correspondientes a los investigadores que conformamos el grupo.

Agradeciendo de antemano su valiosa atención y la receptividad brindada a nuestras iniciativas, me despido.

\closing{Atentamente,}

\end{letter}
\end{document}
